Performance for the first variant is shown in \autoref{fig:pdl1ex1_accuracy}.
Performance barely increases as $\nq$ increases, though there appears to be some effect for lower $Q_y$.
Importantly, performance for low $Q_y$ is moderate, while performance for high $Q_y$ is bad.
For $Q_y = 0.9,$ for instance, accuracy is close to $0.9$,
while an accuracy of $0.9$ can be achieved with an empty ruleset.
The same is observed for $Q_y = 0.8$.
For $Q_y = 0.5$, however, an accuracy of over $0.70$ is achieved.

Performance for the second variant is shown in \autoref{fig:pdl1ex1neg_accuracy}.
The analysis is analogous.
Notably, the extra features yield modest increase in performance overall,
suggesting that lower-bounds may be important for finding relevant rulesets.

This method also gives results on the best possible rulesets that exist.
Its conclusion that, in a majority of the cases, there do not exist meaningfully better rulesets,
in the sense that no meaningfully higher accuracy can be reached.
For further details we refer the reader to \autoref{sec:boolean_model_appendix}.

Based on these results, choosing $Q_y = 0.5 $ or $Q_y = 0.6$ is preferred.
Nevertheless, the results are arguably underwhelming,
compared to the accuracies that are achieved on benchmark datasets.
While it may be possible to find consistent rulesets describing relevant phenomena,
rulesets of this form with moderate complexity ($C=20$) are not good predictors for thresholded PD-L1 expression.

Still, we are interested in the behavior of the model as we restrict complexity.
How much does performance drop?
What do simple rulesets look like?
Are they consistent among splits of the dataset?

To study these results, we will tailor the model to thresholds to speed up computation
and to set the scene for future extensions.


\begin{figure}[h]
    \centering
    \includesvg[width=0.8\textwidth]{gen/PDL1ex1/fig_PDL1ex1_accuracy.svg}
    \caption{Ruleset performance for the first variant of the experiment,
        only considering binary features \eqref{eq:BM_TF_activity_quantiles}.}
    \label{fig:pdl1ex1_accuracy}
\end{figure}

%\begin{figure}[h]
%    \centering
%    \includesvg[width=0.8\textwidth]{gen/PDL1ex1/fig_PDL1ex1_solve_time.svg}
%    \caption{Pricing phase solve time for the first variant of the experiment, only considering binary features \eqref{eq:BM_TF_activity_quantiles} }
%    \label{fig:pdl1ex1_solve_time}
%\end{figure}

\begin{figure}[h]
    \centering
    \includesvg[width=0.8\textwidth]{gen/PDL1ex1neg/fig_PDL1ex1neg_accuracy.svg}
    \caption{Ruleset performance for the first variant of the experiment,
        considering both features \eqref{eq:BM_TF_activity_quantiles} and \eqref{eq:BM_TF_activity_quantiles_reverse}.}
    \label{fig:pdl1ex1neg_accuracy}
\end{figure}

%\begin{figure}[h]
%    \centering
%    \includesvg[width=0.8\textwidth]{gen/PDL1ex1neg/fig_PDL1ex1neg_solve_time.svg}
%    \caption{Pricing phase solve time for the first variant of the experiment,
%        considering both features \eqref{eq:BM_TF_activity_quantiles} and \eqref{eq:BM_TF_activity_quantiles_reverse}}
%    \label{fig:pdl1ex1neg_solve_time}
%\end{figure}